\naslov{Vrednotenje napovednih modelov}

Recimo, da obstaja idealen napovedni model $Y = m(\mb{X}) + \varepsilon$, kjer
je za napako $E(\varepsilon) = 0$.
Zanima nas pričakovana napaka napovednega modela $\hat{m}$, naučenega na neki
množici $S$, pri $\mb{X} = \mb{x}_0$.
Model $m$ je determinističen, torej velja
\[
  E(Y \such \mb{X} = \mb{x}_0) = E(m(\mb{x}_0) + \varepsilon) = E(m(\mb{x}_0)) =
  m(\mb{x}_0)
\]
in
\[
  \var(Y \such \mb{X} = \mb{x}_0) = E(\varepsilon^2) = \sigma_\varepsilon^2.
\]
S podobno enostavnim računom pridemo do
\[
  E(\operatorname{Err}(\mb{x}_0))
  = E((Y - \hat{m}(\mb{x}_0))^2)
  = \sigma_\varepsilon^2 + \left( E(\hat{m}(\mb{x}_0)) - m(\mb{x}_0) \right)^2
  + \var(\hat{m}(\mb{x}_0)).
\]
To nas vodi do definicije dveh novih količin: \pojem{pristranskost}
\[
  E(\hat{m}(\mb{x}_0)) - m(\mb{x}_0)
\]
in \pojem{varianca}
\[
  \var(\hat{m}(\mb{x}_0)).
\]

\vprasanje{Izpelji predpis za pričakovano napako napovednega modela. Kaj sta
  pristranskost in varianca?}

Pri učenju modelov ločimo \pojem{učno} in \pojem{testno} napako.
Učna napaka je oblike
\[
  \operatorname{Err}_{\text{train}}(m, S)
  = E(L(y, m(\mb{x})) \such (\mb{x}, y) \in S)
  = \inv{\abs{S}} \sum_{(\mb{x}, y) \in S} L(y, m(\mb{x})),
\]
kjer je $L$ funkcija izgube.
Ta napaka je seveda optimistična, zato merimo tudi testno napako
\[
  \operatorname{Err}_{\text{test}}(m, S)
  = E(L(y, m(\mb{x})) \such (\mb{x}, y) \notin S).
\]
Razliki med tema napakama pravimo \pojem{optimizem}
\[
  o = \operatorname{Err}_{\text{test}} - \operatorname{Err}_{\text{train}}.
\]
Izkaže se, da v splošnem velja
\[
  E(o) = \frac{2}{\abs{S}} \sum_{(\mb{x}, y) \in S} \cov(m(\mb{x}), y),
\]
torej je optimizem obratno sorazmeren številu učnih primerov $\abs{S}$.
V skrajnem primeru za $\abs{S} \to \infty$ sploh ne potrebujemo testnih
primerov.
Pri linearnih modelih se izkaže $\cov(m(\mb{x}), y) = p \sigma_\varepsilon^2$,
torej je optimizem premo sorazmeren s številom napovednih spremenljivk.

\vprasanje{Kaj je optimizem? Kako se izraža v primeru linearnega modela?}

Poznamo več načinov za izbiro učnih in testnih podatkov.
Najbolj enostaven je naključno vzorčenje, kjer vzamemo nekaj naključnih
primerov za učne, ostale pa pustimo za testne.
Problem s tem je, da je ocena napake občutljiva na izbiro primerov.

Boljši način je prečno preverjanje, kjer razdelimo $S$ na $k$ enako velikih
množic, ki imajo približno enako porazdelitev ciljne spremenljivke.
Potem naredimo $k$ modelov, vsakič vzamemo en kos za testno in ostale za učno
množico.
Končna napaka je potem povprečje napak na posamičnih testnih.

\vprasanje{Razloži prečno preverjanje.}

Druga možnost je zankanje, kjer vzamemo $B$ naključnih vzorcev množice $S$ s
ponavljanjem.
Te množice so vse enako velike kot $S$, imajo vlogo učnih, njihov komplement pa
vlogo testnih množic.
Verjetnost, da nek primer ni v učni množici, je potem $(1 -
\inv{\abs{S}})^{\abs{S}} \approx e^{-1}$.
Napako izmerimo za vsak model (na komplementu njegovih učnih podatkov), končna
ocena napake je povprečje teh.

\vprasanje{Razloži zankanje.}

Pri dvojiški klasifikaciji imamo dva pristopa: ali napovedujemo razred, ali
napovedujemo verjetnost.
Če napovedujemo razred, je funkcija izgube $L$ ravno indikator, če smo se
zmotili.
Problem tu je, da ta funkcija izgube ne loči med lažnimi pozitivnimi in lažnimi
negativnimi primeri, kar bomo poskusili popraviti kasneje.
Pred tem si oglejmo \pojem{napovedno napako}
\[
  \text{Err} = \frac{\text{FP} + \text{FN}}{n}
\]
in \pojem{napovedno točnost}
\[
  \text{Acc} = \frac{\text{TP} + \text{TN}}{n} = 1 - \text{Err}.
\]

\vprasanje{Kaj sta napovedna napaka in točnost pri dvojiški klasifikaciji?}

Boljši meri sta \pojem{delež pravilno razvrščenih primerov} ali
\pojem{občutljivost}
\[
  \text{TPR} = \frac{\text{TP}}{\text{P}},
\]
ki ima za idealne modele vrednost $1$, in \pojem{delež napačno razvrščenih
  primerov}
\[
  \text{FPR} = \frac{\text{FP}}{\text{N}} = 1 - \frac{\text{TN}}{\text{N}},
\]
ki ima v idealnem primeru vrednost $0$.
Ti meri lahko uporabimo kot ordinatno in abscisno os v t.i.~prostoru ROC --
\pojem{receiver operating characteristic}.
Modeli z enako napako so v tem prostoru na isti (pozitivni) diagonali, z
idealnim modelom v zgornjem levem kotu in naključnim modelom na simetrali lihih
kvadrantov.
Če je ena napaka tipa FP ali FN slabša kot druga, lahko primerno naklonimo
premice enakovrednih modelov, (ki zdaj ne bodo $y = x + c$, temveč $y = kx +
c$), in izbiramo med modeli na primeren način.

\vprasanje{Kaj je prostor ROC? Kako v njem poiščeš idealni model, če so napake
  FP $\alpha$-krat slabše od napake FN?}

Pri klasifikaciji s pomočjo verjetnosti si zadamo prag $\theta$, da je napoved
enaka $\mathbbm{1}(m(\mb{x}) \ge \theta)$.
Če spreminjamo $\theta$ od $0$ do $1$, se pomikamo v ROC diagramu od levega
spodnjega kota do desnega zgornjega.
Idealen model potem izberemo na podoben način kot prej, s sweepom premice
določenega naklona.

Ploščina pod krivuljo ROC nam poda dodatno mero AUC, ki je neodvisna od izbire
$\theta$ (oz.~se nanaša na vse možne izbire $\theta$).
Poda nam oceno razdalje med verjetnosti za pozitivne in negativne primere.

\vprasanje{Kako izbereš optimalen odločitveni prag pri klasifikaciji z
  verjetnostjo?}

% LocalWords:  lažnimi
